% Experiment plan for the Results section.
% Required packages: amsmath,amssymb,booktabs,tabularx
% This is a protocol/design document. Replace protocol language with observed
% results after the corresponding analyses have been completed.

\section{Results}
\label{sec:results}

\subsection{Feature-MAE learns stable urban visual representations}
\label{sec:results_representation}

\paragraph{Claim and evaluation principle.}
This section will test whether Feature-MAE learns a reproducible representation
of urban scenes rather than merely memorizing the training sample. Evaluation
will be performed on a fixed, panorama-disjoint test set containing 256 clean
images per city ($n=7{,}680$). These images, and all other headings belonging to
the same panoramas, must be excluded before retraining the final models. The
existing 128-image-per-city validation split will remain dedicated to checkpoint
selection. All inferential summaries will treat panoramas, rather than patches,
as the independent sampling unit.

\paragraph{Experiment R1: masked-feature reconstruction.}
For each test image, 75\% of DINOv3 patches will be masked using ten independent
mask realizations. Reconstruction quality will be quantified by masked-patch
cosine similarity,
\begin{equation}
  R_i = \frac{1}{|M_i|}\sum_{p\in M_i}
  \frac{\mathbf{t}_{ip}^{\top}\widehat{\mathbf{t}}_{ip}}
  {\lVert\mathbf{t}_{ip}\rVert_2
   \lVert\widehat{\mathbf{t}}_{ip}\rVert_2},
\end{equation}
and reported as a city-stratified mean with a 95\% confidence interval obtained
by bootstrapping panoramas within cities. We will additionally report the 10th,
50th and 90th percentiles, preventing a small number of easy scenes from
dominating the mean. A spatially permuted control, in which visible patch
positions are shuffled within images, will test whether reconstruction depends
on spatial context rather than only on the marginal distribution of DINOv3
features.

\paragraph{Experiment R2: matched baselines.}
Two classes of baseline will be kept separate because they answer different
questions. Reconstruction baselines will include (i) a global mean-token
predictor, (ii) a position-specific mean-token predictor, and (iii) a dense
autoencoder with the same 512-dimensional width and comparable parameter count
but no masking. Representation baselines will include (iv) a 512-component PCA
projection and (v) the existing Top-$K$ sparse autoencoder with width 512 and
$K=32$. Feature-MAE and the reconstruction baselines will be compared using the
same masked-cosine test. PCA, Top-$K$ SAE and Feature-MAE will be compared on
representation properties that do not privilege a particular training loss:
dead-dimension rate, effective rank, mean absolute inter-dimension correlation,
spatial coherence of activation maps, cross-seed reproducibility and external
morphology localization (Experiment R12). The Top-$K$ model should not be
compared directly with Feature-MAE using masked reconstruction because it was
not trained to infer missing contextual features.

\paragraph{Experiment R3: hyperparameter sensitivity.}
We will vary one design choice at a time around the final configuration:
mask ratio $\rho\in\{0.50,0.75,0.90\}$, latent width
$W\in\{256,512,1024\}$, and training-city count
$C\in\{10,20,30\}$. The 10- and 20-city subsets will be geographically
stratified and fixed before analysis. Screening runs will use a common
city-balanced subset of 3,200 images per city and five epochs; the selected
configuration and its nearest competitors will then be rerun on the full
training set. Outcomes will include held-out reconstruction, effective rank,
dimension redundancy, mean heatmap total variation and matched-dimension
stability. Sensitivity conclusions will be based on effect sizes and bootstrap
confidence intervals, not on the best single run.

\paragraph{Experiment R4: optimization and seed stability.}
The final configuration will be trained from seeds 42, 43 and 44. We will show
training and validation losses for every epoch, the generalization gap, gradient
norm and peak GPU memory. Because latent dimensions can permute across runs,
dimensions will be aligned with the Hungarian algorithm using a similarity
score that averages encoder-direction cosine similarity and activation-profile
correlation on the fixed test set. Stability will be summarized by matched
encoder and decoder cosine similarity, Jaccard overlap of the top-100 responding
panoramas, patch-level heatmap correlation and adjusted Rand index (ARI) of the
resulting 32/64 partitions. A seed-stable representation should show smooth
convergence, a small train--validation gap and substantially stronger matched
agreement than randomly permuted dimensions.

\paragraph{Planned presentation.}
The main figure should contain (a) training and validation curves, (b) held-out
reconstruction by city, (c) the baseline comparison across reconstruction,
compactness and stability metrics, and (d) the cross-seed matched-dimension
distribution. Mask-ratio, width and city-count sweeps should be reported in an
Extended Data figure, with complete numeric results in a supplementary table.

\subsection{Latent visual dimensions form a multiscale hierarchy}
\label{sec:results_hierarchy}

\paragraph{Claim and unit of analysis.}
This section will test whether the 512 learned dimensions form reproducible,
nested families when similarity is defined jointly in parameter, spatial,
geographic and activation-context spaces. The primary hierarchy will remain the
single Ward tree obtained from the 32-dimensional spectral embedding of the
five-view rank-fused affinity matrix. The prespecified tree cuts are 32 coarse
and 64 fine clusters.

\paragraph{Experiment R5: contribution of the five similarity views.}
We will first visualize the encoder, linearized-decoder, patch-position,
city-profile and PPMI activation-context similarity matrices in the common
dendrogram order. Pairwise agreement between views will be measured with a
Mantel correlation using 9,999 permutations and by overlap of each dimension's
30-nearest-neighbour set. We will then perform leave-one-view-out ablations and
five single-view clusterings. Each alternative will be compared with the full
hierarchy using ARI, normalized mutual information (NMI), variation of
information and dimension-wise cluster Jaccard overlap. This analysis will show
whether the hierarchy is supported by multiple sources of evidence or driven by
one dominant view.

\paragraph{Experiment R6: hierarchy and cluster resolution.}
The complete dendrogram will be displayed together with the 32- and 64-cluster
annotation strips. For transparency, we will report cluster-size distributions,
normalized size entropy, silhouette score, Calinski--Harabasz index and
Davies--Bouldin index for cuts from $K=16$ to $K=96$. The purpose of this scan is
diagnostic rather than post-hoc optimization: $K=32$ and $K=64$ remain the
prespecified reporting resolutions because they provide a two-level hierarchy
and because both are obtained from the same tree. The 64-to-32 parent mapping
will be reported in full as a supplementary table.

\paragraph{Experiment R7: cluster stability.}
Three complementary perturbations will be used. First, the existing interleaved
half-sample analysis will be retained. Second, 200 city-stratified panorama
bootstraps will recompute the three activation-derived views, rank fusion,
spectral embedding and Ward tree while keeping learned model parameters fixed.
Third, leave-one-city-out analyses will repeat the hierarchy 30 times to identify
clusters that depend disproportionately on one city. For each perturbation, we
will report per-cluster Jaccard stability, global ARI/NMI and the fraction of
fine clusters retaining the same coarse parent. Null distributions will be
created by independently permuting dimension labels in each similarity view.
Cluster stability will be considered convincing only when observed agreement
exceeds the 95th percentile of the corresponding null distribution.

\paragraph{Experiment R8: seed-to-hierarchy reproducibility.}
For each of the three Feature-MAE seeds in Experiment R4, the full five-view
hierarchy will be rebuilt after Hungarian dimension alignment. We will compare
cophenetic distance matrices and the 32/64 partitions across seeds. This test is
stricter than resampling images with a fixed model because it includes
uncertainty from representation learning itself.

\paragraph{Planned presentation.}
The current publication dendrogram should form the central panel. Supporting
panels should show (a) the five reordered similarity matrices, (b) the
leave-one-view-out ARI/NMI values, and (c) cluster-specific bootstrap and
leave-one-city-out stability. The full resolution scan and seed-level
cophenetic comparisons should be placed in Extended Data.

\subsection{Urban visual patterns vary systematically across cities}
\label{sec:results_cities}

\paragraph{Cluster prevalence.}
Let $w_{ip}$ denote the winning latent dimension for patch $p$ in image $i$,
and let $\mathcal{D}_k$ be the set of dimensions in cluster $k$. We will define
the panorama-level prevalence of cluster $k$ as
\begin{equation}
  z_{ik}=\frac{1}{196}\sum_{p=1}^{196}
  \mathbb{I}\!\left(w_{ip}\in\mathcal{D}_k\right),
  \qquad \sum_k z_{ik}=1.
  \label{eq:cluster_prevalence}
\end{equation}
When two sampled headings belong to the same panorama, their vectors will be
averaged before city-level aggregation. City prevalence will therefore be the
mean of panorama-level compositions, with 95\% intervals obtained from 1,000
within-city panorama bootstraps.

\paragraph{Experiment R9: cross-city differentiation.}
The $30\times64$ city-by-fine-cluster matrix will be displayed after row-wise
centering and variance scaling for visualization. Statistical comparisons will
use a centered log-ratio transform with a small prespecified pseudocount,
reflecting the compositional nature of the data. City distances will be measured
with Aitchison distance and displayed by principal-coordinate analysis and
hierarchical clustering. Jensen--Shannon distance will be reported as a
robustness analysis. For each cluster, city specificity will be summarized by
normalized entropy; low entropy indicates geographic concentration, whereas
high entropy indicates widespread prevalence. Regional separation will be
tested by PERMANOVA with cities as the permutation units. Because the number of
cities is modest, exact effect sizes and permutation confidence intervals will
be emphasized over asymptotic $P$ values.

\paragraph{Experiment R10: within-city geographic organization.}
Panorama-level prevalence will be aggregated to a fixed spatial grid within
each city; grid size will be chosen before examining cluster-specific maps and
kept constant in metric units. For clusters with adequate support, global
Moran's $I$ will quantify spatial autocorrelation using a distance-based
neighbour graph. Significance will be evaluated with 999 coordinate-preserving
permutations within cities and false-discovery-rate correction across clusters.
Maps will be presented for a small, prespecified set of clusters spanning built
form, road space and vegetation, using identical colour limits across cities.
This avoids selecting only visually striking examples after inspecting the
results.

\paragraph{Experiment R11: nuisance robustness.}
We will test whether cross-city differences persist after controlling for
heading, acquisition year and month, image brightness, quality-model score and
the availability of multiple headings per panorama. Four analyses are planned:
(i) recomputation on a heading-balanced subset; (ii) exclusion of the 10\% of
retained images with the highest artefact probability; (iii) adjustment using a
mixed-effects model with city as a random intercept; and (iv) within-panorama
comparison of different headings. Robustness will be summarized by correlations
between original and adjusted city profiles, preservation of city ordination
under Procrustes analysis, and consistency of cluster-specific city ranks. A
negative control will repeat the analysis after shuffling city labels within
acquisition-year strata.

\paragraph{Planned presentation.}
The main figure should combine (a) a city-by-cluster prevalence heatmap, (b) a
composition-aware city ordination, (c) maps for representative clusters in
four to six cities and (d) a coefficient or sensitivity plot showing nuisance
robustness. Complete city profiles, bootstrap intervals and Moran statistics
should be released as source data.

\subsection{Visual patterns correspond to interpretable urban morphology}
\label{sec:results_morphology}

\paragraph{Independent validation strategy.}
Interpretability will be evaluated using information that was not used to train
Feature-MAE or construct the hierarchy. We will combine image-based semantic
masks with geographic road and building attributes. Validation images must be
drawn from the panorama-disjoint test set established for Experiment R1.

\paragraph{Experiment R12: image-level road and building validation.}
An independent pretrained semantic-segmentation model will generate masks for
road, sidewalk, building, vegetation, sky, vehicle and street-furniture classes
on 300 test images per city. For latent dimension $d$ and semantic class $g$, we
will calculate heatmap-weighted enrichment
\begin{equation}
  E_{dg}=
  \frac{\sum_{i,p}A^{+}_{ipd}\,\mathbb{I}(Y_{ip}=g)
        /\sum_{i,p}A^{+}_{ipd}}
       {\sum_{i,p}\mathbb{I}(Y_{ip}=g)/(196N)},
  \qquad A^{+}_{ipd}=\max(A_{ipd},0).
  \label{eq:semantic_enrichment}
\end{equation}
We will report enrichment, patch-level average precision and city-block-bootstrap
confidence intervals. Cluster-level values will be obtained by aggregating
member dimensions without using Qwen labels. A label-free positive result would
be selective enrichment of distinct clusters for road, building and vegetation
rather than uniform responses across all semantic classes.

\paragraph{Experiment R13: external geographic morphology.}
Panoramas will be spatially joined to independently obtained road and building
data. Prespecified variables will include road class, intersection density,
building-footprint coverage, distance to the nearest building edge, local block
size and, where available, building height. Associations with coarse- and
fine-cluster prevalence will be estimated using city-standardized predictors and
mixed-effects models with city-level random intercepts and slopes. Predictive
validity will be assessed by leave-one-city-out cross-validation: models fitted
on 29 cities must predict morphology in the held-out city above permutation
baselines. If complete building data cannot be obtained for all cities, the
confirmatory analysis will use a geographically stratified subset selected
before inspecting cluster associations, and the remaining cities will be
treated as out-of-sample replication sites where data become available.

\paragraph{Experiment R14: VLM-label validation.}
Qwen3-VL-2B will continue to receive exactly ten unique-panorama examples for
each cluster medoid, with original image, heatmap and hotspot crop shown
together. The 64 fine categories will be named before their 32 parents. To
measure reliability, prompts will be repeated three times with independently
ordered evidence sheets. Label consistency will be evaluated using semantic
embedding similarity and agreement in the constrained semantic-type field.
Generic-name rate, parse-failure rate, artefact probability and confidence will
be reported. The VLM outputs will remain strictly post hoc and will not be used
to select the number of clusters or revise membership.

\paragraph{Experiment R15: blinded human audit.}
Three independent raters will review a stratified sample containing all 32
coarse clusters and 32 randomly selected fine clusters. For each category,
raters will see the same ten-image evidence sheet but not the model-generated
name. They will select a semantic type, write a short description, and score
within-category visual coherence on a five-point scale. Agreement will be
summarized with Fleiss' $\kappa$ for semantic type and an intraclass correlation
coefficient for coherence. Qwen labels will be compared with consensus human
descriptions only after consensus has been frozen.

\paragraph{Experiment R16: coarse-to-fine case studies.}
Representative branches will be selected using prespecified criteria: above-
median split-sample stability, at least two fine children and evidence from at
least three cities. For each selected branch, the figure will show the coarse
medoid, its fine-child medoids, positive-activation heatmaps and city prevalence.
Examples should include at least one built-form branch, one road-space branch
and one vegetation or environmental branch. This presentation will demonstrate
how a broad visual family separates into more specific spatial or morphological
patterns without relying on isolated, hand-picked dimensions.

\paragraph{Planned presentation.}
The main figure should contain (a) a cluster-by-external-class enrichment matrix,
(b) cross-validated associations with road and building morphology, (c) VLM and
human agreement, and (d) two or three coarse-to-fine visual branches. Full
evidence sheets, prompts, raw VLM responses and human-rating instructions should
be placed in the Supplementary Information.

\subsubsection*{Statistical safeguards and reporting order}
\label{sec:results_safeguards}

Primary outcomes are panorama-level masked reconstruction, matched-dimension
seed stability, 32/64 cluster stability, city-level compositional prevalence,
and external road/building enrichment. Patch-level observations will never be
treated as independent replicates. Confidence intervals will use panorama
bootstrap within cities or city-block bootstrap, depending on the estimand.
Permutation tests will preserve city and acquisition strata wherever relevant.
False-discovery-rate correction will be applied within each family of
cluster-level tests. Hyperparameters, representative clusters and map colour
limits will be fixed before inspecting the corresponding test results.

The recommended execution order is R1--R4 (representation), R5--R8 (hierarchy),
R9--R11 (cities), and R12--R16 (morphology). This order prevents semantic
interpretation from feeding back into statistical model selection.

\begin{table*}[t]
\centering
\caption{Priority and computational scope of the proposed experiments.}
\label{tab:experiment_priority}
\small
\begin{tabularx}{\textwidth}{@{}lllX@{}}
\toprule
Priority & Experiments & New model training & Purpose \\
\midrule
P0 & R1, R2, R4 & Three full Feature-MAE seeds; matched baselines & Establish reconstruction, baseline performance and reproducibility. \\
P0 & R5--R7 & None & Validate the five-view hierarchy using ablations and resampling. \\
P0 & R9--R11 & None & Quantify city prevalence, spatial structure and nuisance robustness. \\
P0 & R12, R14, R16 & Independent segmentation and repeated VLM inference & Establish interpretable morphology and coarse-to-fine evidence. \\
P1 & R3 & Reduced-data sensitivity grid plus selected full runs & Demonstrate robustness to mask ratio, width and city coverage. \\
P1 & R8 & Uses the three P0 seeds & Propagate training uncertainty into the hierarchy. \\
P1 & R13, R15 & GIS preparation and human annotation & Provide strong external and human validation. \\
\bottomrule
\end{tabularx}
\end{table*}
